% =====================================================================
%  Chương 1 — GIỚI THIỆU TỔNG QUAN
% =====================================================================
\chapter{Giới thiệu tổng quan}
\label{chap:gioithieu}

% =====================================================================
\section{Giới thiệu}
\label{sec:gt_gioithieu}

Trong bối cảnh kinh tế chia sẻ và du lịch số phát triển mạnh, các nền tảng đặt phòng
trực tuyến (Booking.com, Agoda, TripAdvisor, \dots) cùng mạng xã hội du lịch tạo ra
lượng lớn \textbf{đánh giá của khách hàng} (customer reviews) mỗi ngày. Một khách sạn
trung bình có thể nhận hàng trăm đến hàng nghìn phản hồi trong tháng, mỗi phản hồi chứa
thông tin đa chiều về chất lượng dịch vụ, cơ sở vật chất, trải nghiệm và thương hiệu.
Việc đọc và tổng hợp thủ công các đánh giá này không còn khả thi; doanh nghiệp cần
công cụ tự động để trích xuất \emph{ý kiến cụ thể} gắn với từng \emph{khía cạnh}
(aspect) của dịch vụ.

\textbf{Phân tích cảm xúc} (Sentiment Analysis) truyền thống gán một nhãn cảm xúc duy nhất
cho toàn bộ câu hoặc đoạn văn. Cách tiếp cận này thất bại trước các câu chứa nhiều khía
cạnh với cảm xúc mâu thuẫn. Ví dụ:

\begin{quote}
\textit{``The room was beautiful but the staff was rude.''}
\end{quote}

\noindent câu trên thể hiện cảm xúc \textbf{tích cực} đối với \textit{room} (FACILITY)
và \textbf{tiêu cực} đối với \textit{staff} (SERVICE). \textbf{Phân tích cảm xúc theo
khía cạnh} (Aspect-Based Sentiment Analysis -- ABSA) \cite{pontiki2016semeval} giải quyết
hạn chế này bằng cách phân tích cảm xúc \emph{theo từng khía cạnh} xuất hiện trong văn bản.

Các mô hình ngôn ngữ tiền huấn luyện dựa trên kiến trúc \textbf{Transformer}
\cite{vaswani2023attentionneed}, đặc biệt BERT \cite{devlin2019bert} và T5
\cite{raffel2020t5}, đã đạt kết quả state-of-the-art trên nhiều bài toán ABSA.
Kiến trúc \textbf{FAST-LCF-BERT} \cite{zeng2019lcf, yang2021lcfatepc} kết hợp cơ chế
\textbf{Local Context Focus (LCF)} -- tập trung biểu diễn vào vùng ngữ cảnh quanh khía
cạnh -- giúp cải thiện đáng kể độ chính xác phân loại cảm xúc (Aspect Polarity
Classification, APC).

Tuy nhiên, lợi ích về độ chính xác đi kèm \textbf{chi phí tính toán cao}: cơ chế
self-attention có độ phức tạp $O(L^2)$ theo độ dài chuỗi token $L$. Trong triển khai
thực tế -- ví dụ trung tâm chăm sóc khách hàng cần xử lý hàng nghìn đánh giá mỗi ngày
-- chi phí này trở thành rào cản. \textbf{Token Merging (ToMe)} \cite{bolya2023tome},
ban đầu đề xuất cho Vision Transformer, gộp các token có biểu diễn tương tự nhau để
rút ngắn chuỗi mà không cần huấn luyện lại từ đầu.

Khóa luận \textit{Attention-Based Token Merging for Aspect-Based Sentiment Analysis} tập
trung vào câu hỏi: \emph{khi kết hợp Token Merging với LCF trong mô hình ABSA đa nhiệm,
chiến lược gộp token nào bảo toàn thông tin tốt nhất cho phân loại cảm xúc và nhóm khía
cạnh?} Hệ thống được hiện thực hoàn chỉnh trong repository \texttt{thesis\_apc\_baseline},
bao gồm mã nguồn huấn luyện, đánh giá và giao diện demo.

\subsection{Bối cảnh ứng dụng trong ngành khách sạn}

Trong ngành lưu trú, phản hồi khách hàng thường đề cập đồng thời nhiều khía cạnh:
\texttt{SERVICE} (nhân viên, lễ tân), \texttt{FACILITY} (phòng, hồ bơi), \texttt{AMENITY}
(tiện nghi), \texttt{EXPERIENCE} (trải nghiệm tổng thể), \texttt{BRANDING} (thương hiệu)
và \texttt{LOYALTY} (chương trình khách hàng thân thiết). Một hệ thống ABSA đầu-cuối cần:

\begin{enumerate}
  \item \textbf{Trích xuất} cụm từ khía cạnh (Aspect Term Extraction -- ATE) xuất hiện
        trong câu;
  \item \textbf{Phân loại} đồng thời cảm xúc (Positive/Negative/Neutral) và nhóm khía
        cạnh (APC) cho mỗi cụm từ;
  \item \textbf{Tổng hợp} thành bộ ba \((a, c, s)\) -- đầu ra cuối cùng phục vụ báo cáo
        và phân tích kinh doanh.
\end{enumerate}

Khóa luận hình thức hóa bài toán trên thành \textbf{Joint Triplet Extraction} trên bộ
dữ liệu đánh giá khách sạn định dạng \texttt{.apc} (mục~\ref{sec:kq_data},
Chương~\ref{chap:ketqua}).

\subsection{Động lực nghiên cứu}

Ba động lực chính dẫn đến đề tài:

\paragraph{(1) Lan truyền lỗi trong kiến trúc pipeline.}
Khi ATE và APC được tách rời (Pipeline), lỗi trích xuất khía cạnh lan sang bước phân loại.
Thực nghiệm trong repo cho thấy ATE F1 giảm từ $79.87$ (Joint) xuống $\approx 60$
(Pipeline), kéo Micro F1 bộ ba giảm khoảng 18 điểm (Chương~\ref{chap:ketqua}).

\paragraph{(2) Chi phí tính toán của Transformer.}
Self-attention tốn $O(L^2)$; nhiều token (stopword, dấu câu, subword lặp) mang ít thông
tin phân biệt cảm xúc. Token Merging hứa hẹn giảm độ dài chuỗi, nhưng áp dụng trực tiếp
cho ABSA có rủi ro gộp nhầm token thuộc cụm khía cạnh.

\paragraph{(3) Thiếu so sánh hệ thống giữa các chiến lược gộp token khi kết hợp LCF.}
Ngoài chiến lược bipartite gốc \cite{bolya2023tome}, khóa luận đề xuất và đánh giá thêm
\texttt{sequential\_local} (quét lân cận, bảo vệ ranh giới LCF) và
\texttt{attention\_weighted} (ưu tiên gộp token ít attention, bảo vệ token khía cạnh).

% =====================================================================
\section{Mục tiêu}
\label{sec:gt_muctieu}

Khóa luận hướng tới các mục tiêu cụ thể sau:

\begin{enumerate}
  \item \textbf{Xây dựng hệ thống ABSA đầu-cuối} trích xuất bộ ba
        \((\text{aspect term}, \text{aspect category}, \text{sentiment})\) từ câu đánh
        giá khách sạn, gồm bốn stage: tiền xử lý, ATE (T5), APC
        (\texttt{FastLcfBertMultiTask}), tổng hợp bộ ba (Chương~\ref{chap:phuongphap}).

  \item \textbf{Tích hợp hai phương pháp kỹ thuật} vào module APC:
        \begin{itemize}
          \item \textbf{Phương pháp 1 (PP1) -- LCF:} tập trung ngữ cảnh cục bộ quanh khía
                cạnh, với hai biến thể CDM (mặt nạ nhị phân) và CDW (trọng số giảm dần);
          \item \textbf{Phương pháp 2 (PP2) -- Token Merging:} gộp token dư thừa sau
                backbone, với ba chiến lược \texttt{bipartite}, \texttt{sequential\_local},
                \texttt{attention\_weighted}.
        \end{itemize}

  \item \textbf{Đánh giá thực nghiệm có hệ thống} 12 cấu hình trên hai backbone
        (\texttt{bert-base-uncased}, \texttt{t5-base}), hai kiến trúc Joint/Pipeline,
        với các độ đo Micro/Macro F1, Oracle Joint Acc và thời gian huấn luyện
        (Chương~\ref{chap:ketqua}).

  \item \textbf{Phân tích} ảnh hưởng của LCF, chiến lược ToMe, mất cân bằng nhãn và lỗi
        ATE đến hiệu suất bộ ba; chỉ ra hướng tăng tốc thực sự (compact, pre-BERT ToMe)
        cho nghiên cứu tiếp theo (Chương~\ref{chap:ketluan}).
\end{enumerate}

% =====================================================================
\section{Câu hỏi nghiên cứu}
\label{sec:gt_cauhoi}

Từ mục tiêu trên, khóa luận đặt ra các câu hỏi nghiên cứu (Research Questions):

\begin{description}
  \item[RQ1.] \textbf{LCF có cải thiện phân loại APC so với baseline không?}
        CDM và CDW ảnh hưởng thế nào đến Micro F1 và Macro F1, đặc biệt trên nhãn thiểu số?

  \item[RQ2.] \textbf{Token Merging có làm suy giảm chất lượng biểu diễn không?}
        Ba chiến lược gộp (\texttt{bipartite}, \texttt{sequential\_local},
        \texttt{attention\_weighted}) so với baseline không dùng ToMe?

  \item[RQ3.] \textbf{Kết hợp LCF + ToMe có hiệu quả hơn từng phương pháp riêng lẻ không?}
        Cấu hình nào đạt cân bằng tốt nhất giữa Micro F1, Macro F1 và thời gian huấn luyện?

  \item[RQ4.] \textbf{Backbone BERT và T5 cho APC có khác biệt đáng kể không?}
        So sánh Oracle Joint Acc, Micro F1 và chi phí huấn luyện trên cùng không gian cấu hình.

  \item[RQ5.] \textbf{Lan truyền lỗi ATE ảnh hưởng thế nào đến hiệu suất Pipeline?}
        Oracle Joint Acc có thay đổi giữa Joint và Pipeline với cùng checkpoint APC không?

  \item[RQ6.] \textbf{(Mở rộng)} Dữ liệu bổ trợ (\texttt{supplement}) và tách câu (UOS)
        có cải thiện recall/Micro F1 không? (Kết quả bổ sung ở Bước 4--5, mục~\ref{sec:eval_step4}--\ref{sec:eval_step5}.)
\end{description}

% =====================================================================
\section{Phạm vi}
\label{sec:gt_phamvi}

\subsection{Phạm vi về bài toán}

Khóa luận tập trung vào hai bài toán con ABSA:

\begin{itemize}
  \item \textbf{ATE} -- trích xuất cụm từ khía cạnh, hiện thực bằng T5 sinh có điều kiện
        (\texttt{T5AspectExtractor}, định dạng GAS);
  \item \textbf{APC} -- phân loại đa nhiệm cảm xúc (3 lớp) và nhóm khía cạnh (6 lớp),
        hiện thực bằng \texttt{FastLcfBertMultiTask} với LCF và ToMe.
\end{itemize}

Các bài toán phức tạp hơn (ASTE, ASQP, trích xuất ý kiến) \textbf{không} nằm trong phạm
vi thực nghiệm trực tiếp, chỉ được khảo sát ở mức lý thuyết (mục~\ref{subsec:absa_subtasks},
Chương~\ref{chap:coso_lythuyet}).

\subsection{Phạm vi về dữ liệu}

\begin{itemize}
  \item \textbf{Miền:} đánh giá khách sạn (hotel reviews), định dạng \texttt{.apc};
  \item \textbf{Quy mô:} train 2\,448 / dev 304 / test 312 mẫu;
  \item \textbf{Ngôn ngữ:} dữ liệu trong repo hiện tại là \textbf{tiếng Anh} (các câu
        đánh giá khách sạn quốc tế); README mô tả bối cảnh tiếng Việt -- người đọc cần
        thống nhất khi báo cáo;
  \item \textbf{Nhãn:} 6 category $\times$ 3 sentiment; đánh giá bộ ba yêu cầu khớp cả
        ba thành phần.
\end{itemize}

\subsection{Phạm vi về mô hình và thực nghiệm}

\begin{itemize}
  \item Backbone APC: \texttt{bert-base-uncased} và \texttt{t5-base} (encoder);
  \item ATE: \texttt{t5-base}, beam search $=4$;
  \item ToMe: \textbf{chỉ dùng chế độ \texttt{resize=True}} trong thực nghiệm chính -- đo
        \emph{chất lượng biểu diễn}, không đo tăng tốc suy luận (mục~\ref{subsubsec:why_resize},
        Chương~\ref{chap:phuongphap});
  \item Huấn luyện: \texttt{clause\_split\_mode=none}; \texttt{USE\_SUPPLEMENT=False} trong
        script mặc định \texttt{run\_joint\_experiments.py} (thí nghiệm supplement ở Bước 4);
  \item Môi trường: Python $\ge$3.9, PyTorch $\ge$2.0, CUDA khuyến nghị; seed $=42$.
\end{itemize}

\subsection{Phạm vi ngoài đề tài}

Các thành phần sau có trong repo nhưng \textbf{không} là trọng tâm chính của luận văn:

\begin{itemize}
  \item \textbf{GAS} (\texttt{gas/}): T5 sinh bộ ba một bước -- so sánh paper ở Bước 3;
  \item \textbf{UOS} (\texttt{uos/}): tách câu bằng LLM (Qwen 8B qua Ollama) -- Bước 5;
  \item \textbf{Web demo} (\texttt{server/}, \texttt{frontend/}): giao diện trình diễn.
\end{itemize}

% =====================================================================
\section{Cấu trúc luận văn}
\label{sec:gt_cautruc}

Luận văn gồm năm chương chính và phụ lục:

\begin{itemize}
  \item \textbf{Chương 1} (chương hiện tại): giới thiệu, mục tiêu, câu hỏi nghiên cứu, phạm vi;
  \item \textbf{Chương 2}: cơ sở lý thuyết -- ABSA, Transformer, BERT/T5, LCF, ToMe;
  \item \textbf{Chương 3}: phương pháp thực hiện -- kiến trúc bốn stage, PP1/PP2, ví dụ minh họa;
  \item \textbf{Chương 4}: kết quả thực nghiệm -- năm bước đánh giá, phân tích, error analysis;
  \item \textbf{Chương 5}: kết luận, hạn chế, hướng phát triển;
  \item \textbf{Phụ lục}: cấu trúc mã nguồn, hướng dẫn tạo hình, bảng F1 theo nhãn.
\end{itemize}

Mã nguồn và kết quả thực nghiệm nằm tại repository \texttt{thesis\_apc\_baseline}; các
file CSV trong \texttt{runs\_ate/} là nguồn số liệu cho các bảng ở Chương~\ref{chap:ketqua}.

% =====================================================================
\section{Tổng quan đóng góp}
\label{sec:gt_donggop}

Đóng góp chính của khóa luận:

\begin{enumerate}
  \item Hiện thực hoàn chỉnh pipeline ABSA bốn stage với mã nguồn tái lập được, gồm 12
        cấu hình thực nghiệm LCF $\times$ ToMe;
  \item Ba chiến lược Token Merging cho chuỗi 1D (văn bản), tích hợp bảo vệ token khía
        cạnh qua \texttt{lcf\_vec};
  \item So sánh hệ thống Joint vs Pipeline, BERT vs T5, CDM vs CDW trên cùng bộ dữ liệu;
  \item Phân tích lỗi chi tiết (ATE boundary errors, nhãn thiểu số F1$=0$) và lập luận
        rõ ràng về chế độ \texttt{resize} vs \texttt{compact}.
\end{enumerate}

% =====================================================================
\section*{Tóm tắt chương}
\addcontentsline{toc}{section}{Tóm tắt chương}

Chương 1 đã trình bày bối cảnh ABSA trong ngành khách sạn, động lực nghiên cứu (lan
truyền lỗi pipeline, chi phí Transformer, thiếu so sánh chiến lược ToMe+LCF), mục tiêu,
câu hỏi nghiên cứu và phạm vi. Chương 2 trình bày nền tảng lý thuyết; Chương 3 mô tả
phương pháp và kiến trúc; Chương 4 báo cáo kết quả thực nghiệm theo năm bước đánh giá.
